% Academic talk (10 slides): overlap-dependent CG iteration growth in ASM,
% from the MFEM<->PETSc bit-identical benchmark through the weighted-ASM ablation.
% Minimal style, complete tables, every data table fully parameter-annotated.
% Compile with XeLaTeX (ctex/Fandol).
\PassOptionsToPackage{table}{xcolor}
\documentclass[aspectratio=169,10pt]{beamer}

\usepackage[fontset=fandol]{ctex}
\usepackage{amsmath,amssymb}
\usepackage{booktabs}
\usepackage{array}

\definecolor{Struct}{HTML}{1F3A5F}
\definecolor{Gray}{HTML}{6E6E6E}

\usetheme{default}
\usefonttheme{professionalfonts}
\setbeamertemplate{navigation symbols}{}
\setbeamercolor{structure}{fg=Struct}
\setbeamercolor{frametitle}{fg=Struct}
\setbeamerfont{frametitle}{series=\bfseries,size=\large}
\setbeamercolor{normal text}{fg=black}
\setbeamertemplate{itemize item}{\textbullet}
\setbeamertemplate{itemize subitem}{--}
\setbeamertemplate{blocks}[default]
\setbeamercolor{block title}{fg=white,bg=Struct}
\setbeamercolor{block body}{bg=Struct!6}
\setbeamertemplate{frametitle}{\vskip0.5em\hspace{0.2em}\usebeamerfont{frametitle}\insertframetitle%
  {\ifx\insertframesubtitle\empty\else\\[0.1em]\small\color{Gray}\insertframesubtitle\fi}\par}
\setbeamertemplate{footline}{\hfill\footnotesize\color{Gray}\ifnum\insertframenumber>12 附录\else\insertframenumber\,/\,12\fi\hspace{0.4cm}\vspace{3pt}}

\newcommand{\ua}{$\uparrow$}
\newcommand{\da}{$\downarrow$}
\newcommand{\yes}{\ensuremath{\checkmark}}
\newcommand{\no}{\ensuremath{\times}}
\renewcommand{\arraystretch}{1.15}
% standard fixed-parameter context line (Sys3), printed under data tables
\newcommand{\ctx}{{\footnotesize\color{Gray}固定参数:Laplace($1$面 Dir$+5$面 Neu,$f{=}1$)· MFEM P1 四面体 $n_x{=}48$($117\,649$ DOF)· $4$ ranks METIS(mult $\in[1,3]$)· rtol$=10^{-6}$,max\_it$=2000$。}}

\begin{document}

% ===== S1 : problem statement =====
\begin{frame}{问题}
{\color{Struct}\bfseries 重叠型加性 Schwarz $+$ 子域 ICC 预条件下,CG 迭代次数随 overlap 增大而\emph{增加}:现象、成因与单层修复\par}
\vskip0.5em
\textbf{主测问题(Sys3,贯穿全部结果).}
\[
-\Delta u = 1\ \ \text{于}\ [0,1]^3,\qquad u=0\ (x{=}0\ \text{面}),\qquad \partial u/\partial n=0\ (\text{其余 }5\ \text{面}).
\]
\textbf{求解器栈.} \ 外层 CG $+$ \texttt{PCASM(BASIC)},overlap $O\in\{0,1,2\}$;子域不完全 Cholesky ICC($L$),$L\in\{0,1,2\}$,\texttt{preonly}。
\vskip0.4em
\textbf{反常.} \ 经典重叠 Schwarz 理论(子域精确解)预期 overlap 增大$\Rightarrow$条件数减小$\Rightarrow$迭代\emph{减少};实测\textbf{相反}。
\vskip0.5em\hrule height0.3pt\vskip0.4em
{\footnotesize 另两个问题作普遍性检验(同样现象,程度不同):\textbf{Sys1} reaction--diffusion $\tfrac1{\Delta t}M+\tfrac12K$(质量主导,良态);\textbf{Sys2} 全 Neumann 椭圆($K$ 奇异,$\ker=\mathrm{span}\{1\}$)。本报告主线为 Sys3。\\[0.2em]
代码与数据:\texttt{github.com/tianhaomahpc-ops/cg-asm-icc-research}。}
\end{frame}

% ===== S2 : MFEM <-> PETSc benchmark =====
\begin{frame}{基准:MFEM 与纯 PETSc 两套实现逐位一致}
\framesubtitle{先确立"数字可信",以便把后续反常归因于数学而非实现}
\begin{columns}[T]
\begin{column}{0.5\textwidth}
\textbf{两套独立实现.}
\begin{itemize}\itemsep0.25em\footnotesize
  \item MFEM 装配矩阵 $\to$ PETSc 求解(\texttt{asm\_demo});
  \item 纯 PETSc \emph{加载同一矩阵} $\to$ 求解(\texttt{pure\_petsc\_load})。
\end{itemize}
\vskip0.3em
\textbf{逐位一致的三层对齐.}
\begin{enumerate}\itemsep0.25em\footnotesize
  \item 矩阵零容差过滤 $|a_{ij}|<10^{-12}$;
  \item 共享同一 METIS 分区;
  \item \texttt{MatLoad} 前用 \texttt{MatSetSizes} 预置逐进程行布局(绕开默认 re-chunk)。
\end{enumerate}
\end{column}
\begin{column}{0.5\textwidth}
\textbf{结果:EQ 的含义.} \footnotesize 同一组参数下,MFEM 与 PETSc 解得的迭代次数\emph{逐位相同}。
\vskip0.3em
\begin{center}\footnotesize
\begin{tabular}{l c c}
\toprule
方案集合 & 组合数 & 一致 \\
\midrule
原方案 $\{0,3,4\}\times O\times L$ & 27 & $27/27$ \\
新方案 $5,6$(全 $O,L,$Krylov) & 15 & $15/15$ \\
\bottomrule
\end{tabular}
\end{center}
\vskip0.3em
\footnotesize \texttt{-ksp\_monitor} 下每步预条件残差范数亦逐位相同。\par
\vskip0.3em
$\Rightarrow$ 后续任何反常都是\textbf{方法的数学性质},非某一边的实现产物。
\end{column}
\end{columns}
\vskip0.5em
\ctx
\end{frame}

% ===== S3 : the anomaly + the flip (wide sweep) =====
\begin{frame}{反常现象与翻转:趋势随 ICC 精度改变}
\framesubtitle{方案 BASIC,外层 CG;完整扫描 $O=0\text{–}5$,$L=0\text{–}6$(迭代次数)}
\begin{center}\footnotesize\setlength{\tabcolsep}{5pt}
\begin{tabular}{c ccccccc}
\toprule
$O\backslash L$ & 0 & 1 & 2 & 3 & 4 & 5 & 6 \\
\midrule
0 & 132 & 102 & 87 & 77 & 72 & 66 & 64 \\
1 & 148 & 102 & 92 & 75 & 65 & 56 & 50 \\
2 & 179 & 117 & 90 & 74 & 65 & 57 & 51 \\
3 & 199 & 130 & 103 & 80 & \textbf{59} & 56 & 50 \\
4 & 203 & 133 & 105 & 83 & 66 & 55 & \textbf{46} \\
5 & 202 & 133 & 106 & 83 & 68 & 56 & 48 \\
\midrule
随 $O$ & \ua & \ua & \ua & $\approx$ & \da & \da & \da \\
\bottomrule
\end{tabular}
\end{center}
\vskip0.2em
\begin{itemize}\itemsep0.2em\footnotesize
  \item 低 ICC 填充($L\le3$):overlap 增大 $\Rightarrow$ 迭代\textbf{上升}(反常);$L{=}0$ 最典型 $132\to202$。
  \item \textbf{翻转点 $L\approx4\text{–}5$}:$L{=}4$ 起 overlap 转为\textbf{有益}($O{=}3$ 取 $59<72$),$L{=}6$ 单调 $64\to46$。
  \item 与成因一致:ICC 越精确 $\Rightarrow$ 因素 A($\omega$)越小 $\Rightarrow$ $C_0^2(\delta)$ 的下降占上风。
\end{itemize}
\ctx
\end{frame}

% ===== S4 : two coupled causes + nail-down table =====
\begin{frame}{成因:两个耦合因素,用"换精确子域解"坐实}
\framesubtitle{方案 BASIC,外层 CG,$L{=}0$;仅改变子域求解器,其余不变}
\textbf{机制.}\ 预条件 CG:$n_{\text{iter}}\sim\tfrac12\sqrt{\kappa}\,\ln(2/\varepsilon)$。抽象加性 Schwarz 上界(Gander, Halpern \& Santugini 2015,定理 2.7;着色形式):
\[
\kappa(P_{\text{ad}})\ \le\
\underbrace{C_0^{2}}_{\substack{\text{稳定分解}\\ \delta\uparrow\Rightarrow\downarrow}}\ \cdot\
\underbrace{\omega}_{\substack{\text{(A) 局部稳定性}\\ \text{精确}{=}1,\ \text{ICC}{>}1}}\ \cdot\
\underbrace{(N_c+1)}_{\substack{\text{(B) 着色数}+\text{粗空间}\\ \delta\uparrow\Rightarrow\uparrow}}
\]
\vskip-0.3em
{\footnotesize 由 $\lambda_{\min}\!\ge\! C_0^{-2}$(稳定分解)与 $\lambda_{\max}\!\le\!\omega(N_c{+}1)$(着色 $+$ 局部稳定性)相除得到;推导见附录。$+1$ 为粗空间 $P_0$,\textbf{本文单层无粗空间} $\Rightarrow C_0^2\omega N_c$。精确子域解 $\omega{=}1$。谱半径形式 $\omega(\rho(\mathcal E){+}1)$ 即该文定理 2.5($=$ Toselli--Widlund 定理 2.7)。}
\begin{center}
\begin{tabular}{l ccc c}
\toprule
子域求解器 & $O=0$ & $O=1$ & $O=2$ & 趋势 \\
\midrule
精确 Cholesky(去掉 A) & 52 & 34 & 30 & \da\ 下降 \\
ICC$(0)$(不精确) & 132 & 148 & 179 & \ua\ 上升 \\
\bottomrule
\end{tabular}
\end{center}
\vskip0.4em
\textbf{结论.}\ 去掉因素 A(子域改精确解),overlap 立即恢复"越大越好"。故反常 $=$ A(不精确 ICC)$\times$ B(BASIC 重复计数)\emph{耦合}所致。\quad\ctx
\end{frame}

% ===== S5 : sASM math + impl + fix =====
\begin{frame}{修复:对称缩放加性 Schwarz(sASM)}
\framesubtitle{去掉因素 B、保持对称、保持单层、保持外层 CG}
\begin{columns}[T]
\begin{column}{0.46\textwidth}
\textbf{数学.}\ 记 $D=\operatorname{diag}(m_k)$,$m_k=$ 自由度 $k$ 所属子域数:
\[
M^{-1}_{\text{sASM}}=D^{-1/2}\,M^{-1}_{\text{BASIC}}\,D^{-1/2}.
\]
{\footnotesize 重叠区的 $m_k$ 重复被两侧 $1/\sqrt{m_k}$ 抵消;$D$ 对角$\Rightarrow$算子对称$\Rightarrow$CG 合法;$O{=}0$ 时 $D{=}I$ 退化为 BASIC。}
\vskip0.3em
\textbf{实现(PCSHELL,3 行 apply).}
{\footnotesize\[ z=D^{-1/2}\!\odot\big(M^{-1}_{\text{BASIC}}\,(D^{-1/2}\!\odot r)\big) \]}
{\footnotesize $D$ 由 \texttt{PCASMGetLocalSubdomains} 的指示向量累加得到。}
\end{column}
\begin{column}{0.54\textwidth}
\textbf{修复结果(外层 CG,$L{=}0$;迭代 / solve-only 秒).}
\vskip0.2em
\begin{center}\footnotesize
\begin{tabular}{l cc cc cc}
\toprule
 & \multicolumn{2}{c}{$O=0$} & \multicolumn{2}{c}{$O=1$} & \multicolumn{2}{c}{$O=2$} \\
\cmidrule(lr){2-3}\cmidrule(lr){4-5}\cmidrule(lr){6-7}
方案 & it & s & it & s & it & s \\
\midrule
BASIC & 132 & 0.19 & 148 & 0.24 & 179 & 0.29 \\
sASM  & 132 & 0.20 & 104 & 0.17 & 103 & 0.17 \\
\bottomrule
\end{tabular}
\end{center}
\vskip0.2em
{\footnotesize $132\to104\to103$,overlap 恢复有效;$O{=}2$ 处迭代与时间均近乎减半。时间为 solve-only(warm,best-of-5),负载机器仅供参考。}
\end{column}
\end{columns}
\vskip0.4em
\ctx
\end{frame}

% ===== S5b : ASM vs sASM, wide sweep + optimal iteration/time =====
\begin{frame}{ASM vs sASM:更大范围扫描下的迭代与时间最优}
\framesubtitle{Sys3,$n_x{=}48$,4 ranks;时间为 solve-only(warm,best-of-5)}
\begin{columns}[T]
\begin{column}{0.5\textwidth}
\textbf{每个 overlap 的最优(取最佳 $L$).}
\vskip0.2em
\begin{center}\footnotesize\setlength{\tabcolsep}{4pt}
\begin{tabular}{c cc cc}
\toprule
 & \multicolumn{2}{c}{BASIC} & \multicolumn{2}{c}{sASM} \\
\cmidrule(lr){2-3}\cmidrule(lr){4-5}
$O$ & it & s & it & s \\
\midrule
0 & 87 & 0.145 & 87 & 0.150 \\
1 & 102 & 0.169 & 75 & 0.127 \\
2 & 90 & 0.175 & \textbf{73} & \textbf{0.124} \\
3 & — & — & 66 & 0.132 \\
\midrule
随 $O$ & \multicolumn{2}{c}{\ua\ 越慢} & \multicolumn{2}{c}{\da\ $O{=}2$ 触底} \\
\bottomrule
\end{tabular}
\end{center}
{\footnotesize BASIC 最优被逼到 $O{=}0$(overlap 纯负担);sASM 把 overlap 变成净收益。}
\end{column}
\begin{column}{0.5\textwidth}
\textbf{全局最优.}
\vskip0.2em
\begin{center}\footnotesize
\begin{tabular}{l c c c}
\toprule
方法 & 配置 & iter & solve-only \\
\midrule
BASIC & $O0,L2$ & 87 & 0.145 s \\
\textbf{sASM} & $O2,L1$ & \textbf{73} & \textbf{0.124 s} \\
\bottomrule
\end{tabular}
\end{center}
\vskip0.3em
\textbf{更大范围(迭代次数).}
\begin{itemize}\itemsep0.18em\footnotesize
  \item sASM 在\textbf{每个 $L$}(含 $L{=}0$)overlap 都有益;BASIC 要 $L\ge4\text{–}5$ 才翻转。
  \item sASM 迭代随 $O,L$ 增大而\textbf{饱和}:$O\ge5$ 后无增益,$L\approx12$ 触底 $\approx\textbf{28}$。
  \item 真下界(精确 Cholesky 子域解)更低($O{=}8$ 达 18),但慢 $20\text{–}40\times$。
\end{itemize}
\end{column}
\end{columns}
\vskip0.2em
\textbf{结论.}\ sASM 最优 \textbf{0.124 s / 73 步},比 BASIC 最优 \textbf{0.145 s / 87 步} 快约 \textbf{14\%} 且迭代更少;关键在 sASM 能用 overlap、BASIC 不能。
\end{frame}

% ===== S6 : unified framework table =====
\begin{frame}{统一框架:对角加权 ASM 变体(后三页的导航图)}
\framesubtitle{两条轴:归一化"剂量"(纵)与"对称与否"(横);sASM 落在"剂量恰好 $+$ 对称"}
\begin{center}
\begin{tabular}{l l c c c c}
\toprule
记号 & $M^{-1}$ & 归一化总量 & 对称 & 适配 Krylov & scheme \\
\midrule
BASIC & $\sum_i R_i^{\top}\widetilde A_i^{-1}R_i$ & 无 & 是 & CG & 0 \\
\textbf{sASM} & $D^{-1/2}M_{\text B}^{-1}D^{-1/2}$ & $D^{-1}$ & 是 & CG & 3 \\
exp2 & $D^{-1}M_{\text B}^{-1}D^{-1}$ & $D^{-2}$(过量) & 是 & CG & 5 \\
exp3 & $D^{-1}M_{\text B}^{-1}$ & $D^{-1}$ & 否 & GMRES/FCG & 6 \\
RAS & $\sum_i \widetilde R_i^{\top}\widetilde A_i^{-1}R_i$ & $D^{-1}$(0/1) & 否 & GMRES/FCG & 0$+$restrict \\
\bottomrule
\end{tabular}
\end{center}
\vskip0.6em
\begin{itemize}\itemsep0.3em
  \item 要抵消 B(重复计数 $\textstyle\sum_i R_i^{\top}R_i=D$)共需一份 $D^{-1}$ 权重;\emph{怎么放、放多少、配什么 Krylov} 即三条轴。
  \item S7 比\textbf{剂量轴}(exp2),S8 比\textbf{对称轴}(exp1$+$exp3),S9 比\textbf{Krylov 轴}(exp5$+$exp4)。
\end{itemize}
\end{frame}

% ===== S7 : dose axis (exp 2) =====
\begin{frame}{剂量轴(实验 2):$D^{-1/2}$ 是正确剂量,过量次优}
\framesubtitle{外层 CG;BASIC / sASM($D^{-1/2}$)/ exp2($D^{-1}$ 两侧),迭代次数}
\begin{center}
\begin{tabular}{cc ccc}
\toprule
$O$ & $L$ & BASIC & sASM $D^{-1/2}$ & exp2 $D^{-1}$ \\
\midrule
0 & 0 & 132 & 132 & 132 \\
1 & 0 & 148 & \textbf{104} & 123 \\
2 & 0 & 179 & \textbf{103} & 128 \\
\midrule
1 & 1 & 102 & 75 & 92 \\
1 & 2 & 92  & 70 & 80 \\
2 & 2 & 90  & 67 & 81 \\
\bottomrule
\end{tabular}
\end{center}
\vskip0.5em
\begin{itemize}\itemsep0.3em
  \item $O{=}0$ 时 $D{=}I$,三者相同(自检)。
  \item $O\ge1$ 时 exp2 \textbf{恒在 BASIC 与 sASM 之间}:两侧 $D^{-1}$(总量 $D^{-2}$)过度归一化,优于 BASIC 但不如正确剂量 $D^{-1/2}$。
\end{itemize}
\vskip0.3em
\ctx
\end{frame}

% ===== S8 : symmetry axis (exp 1 + 3) =====
\begin{frame}{对称轴(实验 1 $+$ 3):非对称 PC 配 CG 失效}
\framesubtitle{子域 ICC$(0)$;比较外层 CG 与 GMRES;迭代次数(DIV $=$ 达 2000 上限未收敛)}
\begin{center}
\begin{tabular}{c cc c}
\toprule
$O$ & RAS $+$ CG & exp3($D^{-1}$单侧)$+$ CG & 对照:RAS $+$ GMRES \\
\midrule
0 & 132 & 132 & 265 \\
1 & \textbf{2000 DIV} & \textbf{2000 DIV} & 232 \\
2 & 111 & \textbf{2000 DIV} & 136 \\
\bottomrule
\end{tabular}
\end{center}
\vskip0.5em
\begin{itemize}\itemsep0.3em
  \item $O{=}0$ 时 RAS$=$exp3$=$BASIC(对称),CG 正常。
  \item $O\ge1$:RAS$+$CG 飘忽($O{=}1$ 发散、$O{=}2$ 又收敛),exp3$+$CG 稳定发散——\textbf{非对称破坏 CG 的收敛保证}。
  \item 同一非对称 PC 改配 GMRES 即正常收敛(右列)。
\end{itemize}
\vskip0.3em
\ctx
\end{frame}

% ===== S9 : Krylov axis (exp 5 + 4) =====
\begin{frame}{Krylov 轴(实验 5 $+$ 4):FCG 救活非对称;GMRES 高迭代源于 restart}
\framesubtitle{子域 ICC$(0)$,$L{=}0$;迭代次数}
\begin{columns}[T]
\begin{column}{0.5\textwidth}
\textbf{实验 5:RAS $+$ FCG(flexible CG).}
\vskip0.2em
\begin{center}\footnotesize
\begin{tabular}{c ccc}
\toprule
$O$ & RAS$+$FCG & RAS$+$CG & sASM$+$CG \\
\midrule
0 & 132 & 132 & 132 \\
1 & \textbf{97} & 2000 DIV & 104 \\
2 & \textbf{93} & 111 & 103 \\
\bottomrule
\end{tabular}
\end{center}
{\footnotesize \texttt{mmax} 取 $30/60/120/600$($O{=}2$)恒为 $93$:短截断即足。FCG 单步比 CG 贵(存 $\sim$mmax 个方向)。}
\end{column}
\begin{column}{0.5\textwidth}
\textbf{实验 4:RAS $+$ GMRES,restart 扫描($O{=}1$).}
\vskip0.2em
\begin{center}\footnotesize
\begin{tabular}{c ccccc}
\toprule
restart & 30 & 60 & 120 & 240 & 600 \\
\midrule
iter & 232 & 147 & 95 & 95 & 95 \\
\bottomrule
\end{tabular}
\end{center}
{\footnotesize 修正 Gram--Schmidt 与经典 Gram--Schmidt 逐行\emph{完全相同}。满 GMRES(restart $\ge120$)下 RAS 随 $O$:$129\to95\to92$,单调下降。}
\end{column}
\end{columns}
\vskip0.5em
\begin{itemize}\itemsep0.25em\footnotesize
  \item FCG 容忍非对称 RAS($O{=}1$ 处普通 CG 发散,FCG 收敛 $97$),且 \texttt{mmax}$=30$ 即达满 GMRES 水平($97/93\approx95/92$)。
  \item 实验 1/3 中 GMRES "看着差" 主要是 \textbf{restart$=30$ 的假象},非正交化损失;放大 restart 即回落。
\end{itemize}
\vskip0.2em\ctx
\end{frame}

% ===== S9b : four-way performance comparison =====
\begin{frame}{四方性能对比与选型(各自调到最优)}
\framesubtitle{Sys3,$n_x{=}48$,4 ranks;时间为 solve-only(warm,best-of-5),归约由 \texttt{-log\_view}}
\begin{center}\footnotesize\setlength{\tabcolsep}{5pt}
\begin{tabular}{l c c c c c c}
\toprule
方法 & 最优配置 & 外层 iter & solve-only & 全局归约 & Krylov 内存 & 对称/CG \\
\midrule
BASIC & $O0,L2$ & 87 & 0.145 s & 310 & $\sim$6 向量 & 是 \\
sASM & $O2,L1$ & 73 & 0.124 s & 285 & $\sim$6 向量 & 是 \\
\textbf{sASM$+$Cheby($d2$)} & $O2,L1$ & \textbf{48} & \textbf{0.115 s} & 278 & $\sim$10 向量 & \textbf{是} \\
\textbf{RAS$+$GMRES($r60$)} & $O3,L2$ & 51 & \textbf{0.102 s} & \textbf{233} & $\sim$62 向量 & 否 \\
\bottomrule
\end{tabular}
\end{center}
\vskip0.4em
\begin{columns}[T]
\begin{column}{0.5\textwidth}
\begin{itemize}\itemsep0.2em\footnotesize
  \item 时间:RAS$+$GMRES $<$ sASM$+$Cheby $<$ sASM $<$ BASIC。
  \item 外层迭代:\textbf{sASM$+$Cheby 最少(48)}(Chebyshev 更好地逼近子块逆,每步零额外全局归约)。
  \item 经典 GS 把正交化合并成一次 \texttt{VecMDot} $\Rightarrow$ GMRES 归约不比 CG 多。
\end{itemize}
\end{column}
\begin{column}{0.5\textwidth}
\textbf{选型.}
\begin{itemize}\itemsep0.2em\footnotesize
  \item 对称$/$CG$/$低内存(心脏 HPC 约束):\textbf{sASM$+$Chebyshev(d2)}——迭代最少、快过 sASM、内存仍低。
  \item 内存充裕$/$可弃对称:\textbf{RAS$+$GMRES} 最快、归约最少,但内存 $\sim$10$\times$。
  \item 最省事:纯 sASM(无 deg$/$restart 可调)。
\end{itemize}
\end{column}
\end{columns}
\end{frame}

% ===== S10 : conclusion + decision =====
\begin{frame}{结论与决策}
\textbf{成因(为何 overlap 增大、迭代上升).}\ overlap 本应通过 $C_0^2(\delta)\!\downarrow$ 改善条件数,但被
\textbf{A}(不精确 ICC 的 $\omega_{\max}/\omega_{\min}$)$\times$\textbf{B}(BASIC 重叠重复计数 $N_c$)耦合反超 $\Rightarrow$ 净上升;去除任一即恢复下降。
\vskip0.4em
\begin{center}\footnotesize
\begin{tabular}{l c c c c l}
\toprule
方法 & 对称 & CG & GMRES & FCG & overlap 趋势($L{=}0$) \\
\midrule
BASIC & 是 & \yes \emph{反常}\ua & — & — & $132\to148\to179$ \\
sASM $D^{-1/2}$ & 是 & \yes \textbf{最优} & — & — & $132\to104\to103$ \\
exp2 $D^{-1}$ 两侧 & 是 & \yes 次优 & — & — & $132\to123\to128$ \\
exp3 $D^{-1}$ 单侧 & 否 & \no 发散 & \yes(restart 限) & \yes & — \\
RAS & 否 & 飘忽 & \yes 满GMRES $129\to95\to92$ & \yes $97\to93$ & — \\
\bottomrule
\end{tabular}
\end{center}
\vskip0.4em
\textbf{各实验最终结论.}
\begin{enumerate}\itemsep0.2em\footnotesize
  \item 剂量:$D^{-1/2}$ 最优,过量($D^{-2}$)次优,无(BASIC)反常上升。
  \item 对称:对称是 CG 的硬约束;非对称 PC 配 CG 失效(exp3 稳定发散、RAS 飘忽)。
  \item Krylov:非对称可由 FCG 救活(短截断即达满 GMRES);GMRES 高迭代源于 restart 而非正交化。
  \item 选型:对称路线选 \textbf{sASM};非对称路线选 \textbf{RAS$+$FCG}。两者均单层、保持 overlap 增大$\Rightarrow$迭代下降。
\end{enumerate}
\end{frame}

% ===== Appendix : derivation of the condition-number bound =====
\appendix
\begin{frame}{附录:上界 $\kappa\le C_0^2\,\omega\,(N_c{+}1)$ 的推导(Gander et al. 2015,定理 2.7)}
\framesubtitle{$P_{\text{ad}}=\sum_{i=0}^{N} P_i$,$P_i=R_i^{\top}\tilde P_i$,$\tilde a_i(\tilde P_iu,v)=a(u,R_i^{\top}v)$;$i{=}0$ 为粗空间}
\footnotesize
\textbf{假设 2.2,2.4 $+$ 着色.}\
(i) \emph{稳定分解}(2.2):每个 $u$ 有 $u=\sum_iR_i^{\top}u_i$ 且 $\sum_i\tilde a_i(u_i,u_i)\le C_0^2\,a(u,u)$;\
(ii) \emph{局部稳定性}(2.4):$a(R_i^{\top}u_i,R_i^{\top}u_i)\le\omega\,\tilde a_i(u_i,u_i)$(精确局部解 $\Rightarrow\omega{=}1$);\
(iii) \emph{着色}(定义 2.6):用 $N_c$ 色给细子域染色,同色支撑不相交($\Rightarrow$ 同色 $a$-正交)。
\vskip0.5em
\textbf{下界(用 i).}\ 由 $u=\sum_iR_i^{\top}u_i$ 得 $a(u,u)=\sum_i a(R_i^{\top}u_i,u)$;每项用 $\tilde a_i$ 的 Cauchy--Schwarz,注意 $\tilde a_i(\tilde P_iu,\cdot)=a(P_iu,u)$:
\[
a(u,u)\le\Big(\textstyle\sum_i\tilde a_i(u_i,u_i)\Big)^{1/2}\Big(\sum_i a(P_iu,u)\Big)^{1/2}
\le\big(C_0^2\,a(u,u)\big)^{1/2}\big(a(P_{\text{ad}}u,u)\big)^{1/2}
\]
$\Rightarrow\boxed{\lambda_{\min}(P_{\text{ad}})\ge C_0^{-2}}$(T\&W 引理 2.5)。
\vskip0.4em
\textbf{上界(用 ii, iii).}\ 按颜色分组细空间 $\sum_{c=1}^{N_c}Q_c$($Q_c=\sum_{i\in I_c}P_i$),同色 $a$-正交 $+$(ii)$\Rightarrow a(Q_cu,u)\le\omega\,a(u,u)$;粗空间另有 $a(P_0u,u)\le\omega\,a(u,u)$。求和:
\[
a(P_{\text{ad}}u,u)=a(P_0u,u)+\sum_{c=1}^{N_c}a(Q_cu,u)\le(N_c{+}1)\,\omega\,a(u,u)\ \Rightarrow\ \boxed{\lambda_{\max}\le\omega(N_c{+}1)}.
\]
\textbf{合并.}\ $\kappa\le C_0^2\,\omega\,(N_c{+}1)$。\ \textbf{本文单层无粗空间 $P_0$ $\Rightarrow$ 去掉 $+1$,即 $C_0^2\,\omega\,N_c$}。\ 谱半径形式 $\omega(\rho(\mathcal E){+}1)$ 为定理 2.5($=$ T\&W 定理 2.7),$\rho(\mathcal E)\le N_k$、$N_c\le N_k$(备注 2.8)。精确子域解 $\omega{=}1$;ICC $\omega{>}1$——因素 A。
\end{frame}

\end{document}
